\scanpage{129}
\mspage{\handwritten{24 bis}}
\deleted{\deleted{$G_b$, où $b\handwritten{\in\Gamma}(G,O_G)$ est la section définie dans 2.7. (NB on \unsure{décomposera|décomposer} $S$ en le préschéma somme de\deleted{\unsure{s}}x préschémas $S_{n,\handwritten{\unsure{\rho|p}}}$, sur chacun desquels $G$ est de dimension relative $n$ et de rang \deleted{réductif} nilpotent $\handwritten{\rho}$, et on construit $b$ séparément sur chacun des $G\handwritten{\vert}S_{n,\handwritten{\rho}}$). D'autre part, comme $\handwritten{\psi}^{-1}(U)$ et $U$ sont lisses donc plats localement de présentation finie sur $S$, pour vérifier que $\handwritten{\psi}$ induit un isomorphisme $\handwritten{\psi}^{-1}(U)\handwritten{\to}U$, il suffit de le vérifier fibre par fibre \inserted{(SGA I 5.7.)}, ce qui nous ramène au cas où $S$ est le spectre d'un corps, que l'on peut supposer algébriquement clos. Il existe alors un sous-groupe de Cartan $C$ de $G$, et si $N$ est le normalisateur, on sait \deleted{\unsure{que}} par le théorème de conjugaison que $\handwritten{\unsure{\mathscr T|\mathscr C}}$ s'identifie à $G/N$, et le morphisme $\handwritten{\psi}\colon X\handwritten{\to}G$ n'est autre que celui défini dans le N\textsuperscript{o} 2. On conclut alors par 2.6. b). Enfin, \deleted{il est aussi} cela montre \deleted{\unsure{q}}\handwritten{a}ussi que $U$ est le \uline{plus grand} ouvert \inserted{$V$} de $G$ tel que $\handwritten{\psi}^{-1}(\deleted{U}V)\handwritten{\to}\deleted{U}V$ soit un isomorphisme, car \deleted{il suffit de le vérifier encore dans} pour vérifier $V\handwritten{\subset}U$, il suffit de le vérifier fibre par fibre.}

\uline{Corollaire} 3.2. \quad So\deleted{\unsure{it}}us les conditions de 3.1. , soit $g$ une section régulière de $G$, i.e. telle que pour tout $s\handwritten{\in}S$, $g(s)$ soit un point réguli\handwritten{er} de $G_s$. Alors il existe un et \handwritten{un} seul sous-groupe de Cartan $C$ de $G$ \deleted{contenant $g(S)$,} tel que $g$ soit une section de $C$.

En effet, l'hypothèse sur $g$ signifie que $g$ est une section de $U$, et la conclusion qu'il existe une unique section de $X$ qui la relève, ce qui \deleted{\unsure{revien|revient}} n'est qu'une autre façon d'exprimer que $\handwritten{\psi}^{-1}(U)\handwritten{\to}U$ est un isomorphisme.

Notons maintenant que l'ouvert $\handwritten{\psi}^{-1}(U)$ du sous-groupe de Cartan $X$ de $G_{\deleted{\unsure{\mathscr T}}\handwritten{\mathscr C}}$ \margin{\handwritten{indice $\mathscr C$}} n'est autre que l'ouvert de $X$ formé des points de $X$ qui sont réguliers dans $G_{\deleted{\unsure{\mathscr T}}\handwritten{\mathscr C}}$ (par quoi on entend: réguliers dans leur fibre). On}
% archive sheet number: 128
